% ---------------------------------------------------------------
\section{What Provenance Can and Cannot Prove}\label{sec:results}
% ---------------------------------------------------------------

\subsection{Exact decidability}

The definition of assurance quantifies over every finite universe. The following lemma,
proved in Appendix~\ref{app:proofs}, removes that quantifier.

\begin{lemma}[Separability]\label{lem:sep}
Every constraint in $\Phi$, every acceptance condition and every goal is either a
condition on the global Booleans alone, or of the form
$\forall i\in\mathcal I.\ \varphi(v_i,g)$, where $v_i\in\{0,1\}^{11}$ is the membership
vector of item $i$ and $g$ the global Booleans. The all-zero vector satisfies every
per-item constraint. Consequently $\mu\models_{\alpha,\mathit{gen}}G$ holds if and only if
it holds over a universe of one item.
\end{lemma}

A single-item universe has $2^{11}$ membership vectors and at most two values of
$\mathit{gs}$, so each instance $(\alpha,\mu,\mathit{gen},G)$ is decided by exhaustive
enumeration. We decided all $6\times 128\times 2\times 6 = 9{,}216$ instances twice:
(i) with the Alloy~6.2.0 analyser~\cite{jackson2019alloy} at scope~5, which includes all
universes of one to five items; and (ii) with an independent enumeration written from the
prose definitions above rather than from the Alloy source. The two agree on all
9,216 instances. An earlier run at scope~1 agreed with scope~5 on all 7,680 instances of
the five original goals, as Lemma~\ref{lem:sep} predicts. The results below are therefore
theorems about the model, not bounded observations. Whether the model is faithful to real
systems is a separate question, which Section~\ref{sec:eval} addresses empirically.

\begin{proposition}[Monotonicity]\label{prop:mono}
If $\mu\models_{\alpha,\mathit{gen}}G$ and $\mu\subseteq\mu'$, then
$\mu'\models_{\alpha,\mathit{gen}}G$.
\end{proposition}

Proposition~\ref{prop:mono} lets every result be stated by its \emph{minimal sufficient}
mechanism sets. Table~\ref{tab:minimal} lists them. In every one of the 72
(profile, pipeline, goal) cases, the set of mechanism sets that achieve the goal is exactly
the upward closure of the listed minimal sets.

\paragraph{Beyond the six profiles}
The six profiles are illustrative. To check that the theorems below do not depend on that
choice, we decided the model for every subset of the seven adversary capabilities:
$128\times128\times2\times6=196{,}608$ decisions. We used a vectorised re-implementation of
the independent checker, validated first against all 9,216 Alloy results. Achievability is
monotone in the mechanisms and anti-monotone in the adversary capabilities. Every theorem
statement below was checked against every instance in its scope, with no mismatch
(artefact: \texttt{theorem\_lattice\_check.json}). The theorems therefore hold for
correlated compromises as well, such as one party that holds both the witness service and
the attestation keys.

\begin{table*}[t]
\centering\small
\caption{Minimal sufficient mechanism sets (exact for the model; $\beta$ = either binding
variant, F = \textsf{BindFix}, WS = \textsf{WitSrc}, WI = \textsf{WitIn}, AC = \textsf{AttCap},
AT = \textsf{AttTrans}, Det = \textsf{Detector}). ``---'' means no mechanism set achieves the goal;
$\emptyset$ means the goal holds with no mechanism. Each cell gives filter-only / generative
pipelines when they differ.}\label{tab:minimal}
\begin{tabular}{@{}lllllll@{}}
\toprule
Profile & $G_{\mathrm{bind}}$ & $G_{\mathrm{auth}}$ & $G_{\mathrm{noinf}}$ & $G_{\mathrm{content}}$ & $G_{\mathrm{zero}}$ & $G_{\mathrm{compl}}$\\
\midrule
full & $\beta$ & $\beta{+}$WS, $\beta{+}$AC${+}$AT \;/\; $\beta{+}$WS${+}$AT, $\beta{+}$AC${+}$AT & F & $\beta{+}$Det & --- & AC, WI\\
full+wit & $\beta$ & $\beta{+}$AC${+}$AT & F & $\beta{+}$Det & --- & AC\\
full+tee & $\beta$ & $\beta{+}$WS \;/\; --- & F & --- & --- & WI\\
cap-only & AC, $\beta$ & $\emptyset$ & AC, $\beta$ & Det${+}$AC, Det${+}\beta$, Det${+}$WI & --- & AC, $\beta$, WI\\
exec-only & $\beta$ & $\beta$ & F & $\beta{+}$Det & --- & $\emptyset$\\
trans-only & $\emptyset$ & AT, WS \;/\; AT & $\emptyset$ & Det & --- & $\emptyset$\\
\bottomrule
\end{tabular}
\end{table*}

The table packs six findings. We state them as theorems; the proofs (Appendix~\ref{app:proofs})
are short, because Lemma~\ref{lem:sep} reduces each one to reasoning about a single item.

\subsection{Binding is necessary exactly when the executor is untrusted}

\begin{theorem}[Binding]\label{thm:bind}
If $\textsf{AdvExec}\in\alpha$, then $\mu\models G_{\mathrm{bind}}$ iff
$\beta$. If $\textsf{AdvExec}\notin\alpha$, then $G_{\mathrm{bind}}$ holds iff
$\textsf{AdvCap}\notin\alpha$, or $\textsf{AttCap}\in\mu$ and $\tau$, or $\beta$.
\end{theorem}

No combination of witnesses, attestation of earlier stages or detectors replaces binding
against a corrupt executor, because none of them observes what the executor actually
consumed. The converse matters as much in practice. When training runs in-house on a
trusted executor (the cap-only profile), attested capture alone gives $G_{\mathrm{bind}}$.
In that case the expensive part of the stack, proofs of training, is not needed.

\subsection{Lineage requires binding and an authenticated route}

\begin{theorem}[Lineage under the full adversary]\label{thm:auth}
Against the full adversary,
\begin{align*}
\mu\models G_{\mathrm{auth}}\iff{}& \beta\wedge\big[(\textsf{WitSrc}\wedge(\neg\mathit{gen}\vee\textsf{AttTrans}))\\
&\qquad\vee(\textsf{AttCap}\wedge\textsf{AttTrans})\big].
\end{align*}
If witnesses may collude (full+wit), only $\beta\wedge\textsf{AttCap}\wedge\textsf{AttTrans}$
remains. If attestation is broken (full+tee), only $\beta\wedge\textsf{WitSrc}$ remains, and
only for filter-only pipelines: no mechanism set gives lineage for a generative pipeline.
\end{theorem}

There are two routes to lineage, and they rest on disjoint trust anchors: source
witnesses, or trusted execution of capture and transform. Each fails when its own anchor
fails. Generative pipelines need the transform attested even on the witness route, because
a corrupt transform can otherwise declare an injected item to be synthetic.

\begin{corollary}[Label laundering and declared stage types]\label{cor:launder}
Suppose the auditor accepts synthetic-item claims from a pipeline that declares no
generator, i.e.\ the constraint $\neg\mathit{gen}\Rightarrow\hat S=\emptyset$ is dropped. Then
$\beta\wedge\textsf{WitSrc}$ no longer achieves $G_{\mathrm{auth}}$ for filter-only
pipelines under the full adversary.
\end{corollary}

Corollary~\ref{cor:launder} turns a modelling detail into a design requirement. Stage
types (filter, generator, mixer) must be declared and authenticated, and claims inconsistent
with the declared types must be rejected. We found this requirement because the first
version of the model omitted the constraint. The resulting counterexample has a direct
real-world reading: label poisoned web text ``synthetic'' and the source witnesses never
ask about it.

\subsection{Grinding: binding without fixed randomness leaks influence}

\begin{theorem}[Influence]\label{thm:noinf}
If $\textsf{AdvExec}\in\alpha$, then $\mu\models G_{\mathrm{noinf}}$ iff
$\textsf{BindFix}\in\mu$.
\end{theorem}

The model treats the seed as a single bit of ``prover influence''. The channel itself is
known. Tan et al.~\cite{tan2025zkpot} showed that every existing zero-knowledge proof of
training admits a rejection-sampling attack: the prover trains with many seeds and proves
the run whose model behaves as it wants. With 1,000 seeds, up to 13.4\% of a logistic
regression model's confident predictions could be moved toward the opposite class. Their remedy
removes randomness from the statement: the prover shows that the released model lies within
$\epsilon$ of the optimum of an $L_2$-regularised strongly convex loss, which bounds how far any
accepted model can shift a prediction. That route needs a unique optimum, which LLM training
does not have. The next theorem bounds the influence for any binding check, convex or not, in
terms of how many accepted outputs the prover evaluates.

Fix the committed record $C$. Let $\mathcal A_C$ be the set of outputs that the binding
check would accept for $C$. A \emph{candidate} is any element of $\mathcal A_C$ that the
prover actually produces and evaluates, for example $\mathrm{Train}(C,s)$ for a seed $s$
under deterministic training. Let $\Sigma\subseteq\mathcal A_C$ be the set of candidates,
and $T=|\Sigma|$, bounded by the prover's compute. The candidates are generated independently
of the goal $K$, for example from fresh seeds. Only the \emph{choice} among them may depend
on $K$. The goal $K$ is a random variable, for example a pattern of
predictions the prover wants on chosen inputs. The prover releases some $M\in\Sigma$ that
may depend on $K$ and on all candidates. The verifier accepts because $M\in\mathcal A_C$.

\begin{theorem}[Grinding bound]\label{thm:grind}
$I(K;M\mid C,\Sigma)\le\log_2 T$. If $K$ is uniform on $2^k$ target patterns of the
model's behaviour, then $\Pr[\text{the released model realises }K]\le T\cdot 2^{-k}$, so
realising $k$ chosen bits with probability $\ge 1/2$ requires $T\ge 2^{k-1}$ training runs.
Under \textsf{BindFix} with the seed fixed after commitment and exact replay, $T=1$ and
the bound is $2^{-k}$, i.e.\ no influence beyond chance.
\end{theorem}

\paragraph{What $T$ counts}
The bound counts \emph{evaluated accepted outputs}, not seeds as such. Four consequences
follow, and they separate the mathematics from its implementation.
\begin{enumerate}
\item Every free choice that the check accepts enlarges $\mathcal A_C$: initialisation, data
  order~\cite{shumailov2021ordering}, dropout masks, checkpoint selection. Kaizen, for
  example, has the prover commit its own initial weights and does not fix the number of
  proved iterations in advance~\cite{abbaszadeh2024zkpot}. An optimum-vicinity check~\cite{tan2025zkpot} makes
  $\mathcal A_C$ the $\epsilon$-ball around the optimum; the theorem then still holds, and
  their geometric bound limits what each element of the ball can change.
\item \emph{Numerical tolerance.} A proof-of-learning verifier that accepts every trajectory
  within distance $\delta$ of a replay~\cite{jia2021pol,fang2023pol} accepts many outputs
  per seed. The prover can then search inside the tolerance region, for example by
  perturbing an update. $T$ is the number of such accepted outputs it evaluates, which can
  far exceed the number of seeds. The theorem still applies, but $\mathcal A_C$ is no longer
  small.
\item \emph{Hardware nondeterminism.} If replay is not bit-exact, the verifier must use a
  tolerance, and item 2 applies. If replay is bit-exact, as we measured for fine-tuning on
  one machine (Section~\ref{sec:eval}), each seed yields exactly one accepted output.
\item \emph{\textsf{BindFix}.} The verifier fixes the seed and replays exactly, so
  $|\mathcal A_C|=1$. Its strength depends on the commitment it follows: a prover who can
  commit several records and pick one after seeing the beacon grinds over records instead of
  seeds. Non-equivocation of the record store ($P_2$) closes that door. Tan et al.\ make the
same point for Fiat--Shamir seeds derived from a data commitment, whose own randomness can be
resampled~\cite{tan2025zkpot}.
\end{enumerate}

\subsection{Content: provenance cannot certify the absence of poison}

\begin{theorem}[Content]\label{thm:content}
Let $\textsf{AdvSource}\in\alpha$ and $\tau$ hold.
\begin{enumerate}
\item No mechanism set achieves $G_{\mathrm{zero}}$ unless the detector is assumed perfect
  ($M=\emptyset$); with a perfect detector, $\beta\wedge\textsf{Detector}$ achieves it.
\item If $\textsf{AdvExec}\in\alpha$: $\mu\models G_{\mathrm{content}}$ iff $\beta\wedge\textsf{Detector}$.
\item If $\textsf{AdvExec},\textsf{AdvWit}\notin\alpha$: $\mu\models G_{\mathrm{content}}$ iff
  $\textsf{Detector}\in\mu$ and $\mu\models G_{\mathrm{compl}}$.
\end{enumerate}
If $\tau$ fails, no mechanism set achieves $G_{\mathrm{content}}$.
\end{theorem}

Part~1 is the formal version of a claim that is often made informally and is sometimes
denied~\cite{stokes2021prevent}. Provenance establishes \emph{who} supplied an item and
\emph{that} the model was trained on it. An authorised source that supplies poison passes
every provenance check. The strongest claim available is Part~2's conditional one: every
poisoned item that was used was examined by the detector and missed by it. That claim is
worth exactly the detector's recall. Part~3 is a design result for deployments with a
trusted executor: a training-input witness that makes the record complete can stand in for
binding, so the detector's verdict carries over to the model.

\subsection{Completeness, and why trust anchors should differ}

\begin{theorem}[Completeness]\label{thm:compl}
Let $\textsf{AdvCap}\in\alpha$. Then $\mu\models G_{\mathrm{compl}}$ iff
$(\textsf{AttCap}\in\mu\wedge\tau)$, or $(\textsf{WitIn}\in\mu\wedge\textsf{AdvWit}\notin\alpha)$,
or $(\beta\wedge\textsf{AdvExec}\notin\alpha)$.
\end{theorem}

Attested capture and input witnessing are substitutes for completeness, and they fail
under disjoint assumptions. In Table~\ref{tab:minimal} the full+wit row keeps only
\textsf{AttCap} and the full+tee row keeps only \textsf{WitIn}. A deployment that uses both
survives the loss of either anchor, though not of both. This is the formal argument for
\emph{trust diversity}: build the record from mechanisms whose failure modes are independent.

\subsection{Where a ledger fits}

\begin{proposition}[Ledger]\label{prop:ledger}
The six goals are statements about \emph{accepted} executions. A record store that provides
integrity ($P_1$) and non-equivocation ($P_2$) suffices for every result above. A
transparency log with independent co-signing witnesses provides both, under the assumption
that the witnesses do not collude. A replicated ledger additionally provides inclusion
liveness ($P_8$): an honest party's entry is eventually included despite a censoring
operator. $P_8$ affects availability, not any of the six goals.
\end{proposition}

Proposition~\ref{prop:ledger} is argued, not model-checked (the store model in our
artefact is a bounded check only). It is consistent with the targeted experiment in
Section~\ref{sec:eval}: witnesses expose equivocation but cannot see censorship.
The practical consequence is sharp. A blockchain contributes to training-data assurance
only where censorship-resistant inclusion is a requirement, for example when data
contributors must be able to prove that a deletion request was logged. It does not make
records truthful, complete or poison-free. Those properties come from what is attested,
witnessed and bound \emph{before} anything is written to the store.
